\section{CP2Seq: task and dataset}
\label{sec:cp2seq}

\section{Problem formulation}
\label{sec:problem}

\subsection{Flat-folded states}
\label{sec:problem-states}

A flat-folded state is an isometric map of the sheet into the plane together with a layer
ordering. The map is fixed by the crease pattern; the ordering is not. A single crease pattern
therefore admits many folded states that are geometrically identical and distinct as embeddings,
and determining the ordering is NP hard even when a valid mountain-valley assignment is given
\citep{bern-hayes}.

This distinction is the reason the task cannot be reduced to geometry. Over a folding sequence the
crease graph's topology never changes: the same vertices, edges and faces exist at step one and at
step nineteen. What changes is everything else. Vertex coordinates are reflected at every step.
The overlap structure, which faces lie above which, changes at every step and follows from the
geometry. The layer ordering changes at every step and does \emph{not} follow from the geometry;
it is chosen, and it is where the decisions live. The number of layers accumulates, which is the
formal content of ``it gets harder as it goes''.

\subsection{The action space}
\label{sec:problem-actions}

Simple folds come in named variants, and any benchmark in this area has to say which it uses,
because the complexity results differ between them.

\begin{table}[h]
\caption{Named variants of the simple fold \citep{arkin-map}.}
\label{tab:action-space}
\begin{center}
\begin{tabular}{ll}
\multicolumn{1}{c}{\bf MODEL}  &\multicolumn{1}{c}{\bf DEFINITION} \\
\hline \\
One-layer simple fold    & Rotate a single layer $\pm 180^\circ$ about a line \\
Some-layers simple fold  & Rotate a contiguous run of layers \\
All-layers simple fold   & Rotate every layer the line crosses, as sheet metal bends \\
Finite vs infinite line  & Whether the fold line is a full line or a segment \\
\end{tabular}
\end{center}
\end{table}

CP2Seq instantiates the all-layers and some-layers tiers. The all-layers tier is the simpler
object and it buys three properties at once, the third of which is rarely named. Because the whole
stack moves as a rigid body, no two connected pieces of paper ever move relative to each other
except at the fold line itself, where paper is allowed to bend. Self-intersection is impossible,
the layer ordering is forced rather than chosen, and \textbf{tearing is impossible by
construction.}

Moving only some layers brings all three back, and tearing becomes the binding constraint. If a
moving face is joined to a stationary face along a crease, and that crease is not on the fold
line, the fold would rip the sheet. This is decidable only in the coordinates of the original
square, because adjacency is a fact about the original sheet that survives any amount of folding
and cannot be recovered from the current positions of the polygons alone.

An example makes the tier difference concrete, and it needs only two folds. Fold the right half of
a square over to the left across the vertical midline. The result is a two-layer stack whose top
layer is joined to the bottom layer along the entire crease at the midline. Now select the top
layer alone. Folding its upper half down across the horizontal midline would tear the sheet,
because the moving piece is joined to stationary paper along a vertical crease segment that is
perpendicular to the fold line. Folding its left quarter across a vertical line does not tear,
because the moving piece's only join to stationary paper lies on the fold line itself. Same state,
same layer selection, opposite verdicts, and the only difference is the orientation of the line.

------this 3 paragraphs are over emphasizing on whether all layers fold or some layer'fold  it
gonna tear or not. is that the only point in this part????

\subsection{The task}
\label{sec:problem-task}

Given a crease pattern, produce a sequence of simple folds that reproduces it. Generating an
instance costs one forward pass of the engine; solving it is NP hard. That asymmetry is the
paper's foundation.

---------

It is tempting
to say that legal moves are rare and that this is what makes the task hard. Enumeration refutes
it: legal partial folds \emph{grow} with depth, from 22 at two layers to 332 at thirty-eight, six
times the 56 all-layers folds available at the same state. What falls is the hit rate of uniform
random proposal, from 11.9\% to 3.8\% by seventeen layers, because the space being sampled grows
faster than the legal set inside it. That is a fact about a sampler, not about origami, and
conflating the two would put a false statement about branching factor into the paper.

---------this paragraph should be listed properly and give a proper graph to illustrate it!!!!!!!!



\subsection{Why the corpus is synthesised}
\label{sec:cp2seq-why}

Every sample is built rather than collected, and the reason is that the construction supplies the
label. A folding sequence is not something that can be read off a finished crease pattern by
inspection, which is the whole premise of the task; so a corpus of annotated patterns would require
either a solver, which is the thing being benchmarked, or a human expert, which does not scale and
introduces an error rate nobody can measure. Folding forward and recording avoids both. The sequence
is not inferred after the fact, it is what happened, and the pattern is what the folding left behind.

Generation controls reference sequence depth and the permitted action model. These parameters
let us vary the corpus composition; empirical difficulty is assessed through model and search
performance rather than guaranteed by a generation setting.

\subsection{Generation, step by step}
\label{sec:cp2seq-generation}

The generator starts from a unit square and applies randomly chosen legal simple folds, recording
each one, until it reaches a target depth. It then unfolds and reads off the crease pattern.

{\small\begin{verbatim}
GENERATE(depth_target, action_model):
    state <- unit square, one layer
    seq   <- []

    while len(seq) < depth_target:
        cands <- ENUMERATE_LEGAL_FOLDS(state, action_model)   # all legal folds
        if cands is empty:
            return FAILED               # dead end: discard, never repair
        fold  <- RANDOM_CHOICE(cands)
        state <- APPLY(state, fold)         # the APPLY the verifier uses
        seq.append(fold: line, direction, layers, creases, coupling)

    cp     <- UNFOLD(state)      # creases accumulated over the sequence
    cp     <- PLANARISE(cp)      # split edges at crossings; canonical
    frames <- [state after each step]        # retained as steps.fold

    assert REPLAY(cp, seq) == state      # sample rejected if this fails
    return Sample(cp, seq, frames, metrics(cp, seq))
\end{verbatim}
}

Four properties of this procedure matter for the benchmark and each is a deliberate choice.

\textbf{The generator never searches.} It picks uniformly among legal folds and never backtracks. A
dead end discards the sample rather than repairing it. This keeps generation one forward pass, and
it is also why the recorded sequence carries no claim to being minimal
(Section~\ref{sec:scoring-why}).

\textbf{\texttt{ENUMERATE\_LEGAL\_FOLDS} and \texttt{APPLY} are the same functions the environment
exposes to the model.} The generator is not a privileged path through different code; it is the
environment driven by a random policy. A sample is therefore reachable by definition, using exactly
the action space the model is given.

\textbf{Planarisation is canonical.} A crease pattern is stored with every edge split at every
crossing, so two patterns describing the same geometry have the same edge set and can be compared as
multisets without a tolerance.

\textbf{Every sample is replayed before it is kept.} The assertion is not a test that runs
sometimes; a sample that does not reproduce its own state is discarded at generation time.

\subsection{Generation parameters and difficulty}
\label{sec:cp2seq-difficulty}

\textbf{Reference depth} is the number of folds in the recorded solution, not necessarily the
minimum number needed to solve the instance. \textbf{Coupling} records the number of creases
created by each fold and characterizes how an action affects multiple paper regions. Depth defines
the all-layers strata; coupling is a structural descriptor whose relationship with solver
performance requires separate measurement. Action models specify permitted moves, difficulty
strata group instances, and assistance conditions specify the tools and feedback available.

\subsection{What the corpus contains - This part only need to let readers know that , all layers fold(easy, medium, hard levels) + some layers fold}
\label{sec:cp2seq-contents}

other things like generated or verified, we don't need to mention, we only need to put it in
hugging face or appendix

\subsection{Dataset statistics}
\label{sec:dataset-statistics}

The release contains 600 samples: 400 under the all-layers action model and 200 under the
some-layers action model. Table~\ref{tab:dataset-statistics} reports counts and observed reference
depths from the release manifest. These are corpus counts, not evaluation attempts.

\begin{table}[h]
\centering
\caption{Release composition. Fold counts refer to recorded reference sequences.}
\label{tab:dataset-statistics}
\begin{tabular}{llrr}
\hline
Action model & Stratum or group & Samples & Reference folds \\
\hline
All-layers & Easy & 200 & 4--10 \\
All-layers & Medium & 100 & 11--13 \\
All-layers & Hard & 100 & 14--19 \\
Some-layers & Depth 3 & 50 & 3 \\
Some-layers & Depth 4 & 50 & 4 \\
Some-layers & Depth 5 & 50 & 5 \\
Some-layers & Depth 6 & 50 & 6 \\
\hline
Total & & 600 & 3--19 \\
\hline
\end{tabular}
\end{table}

Tool availability and target-feedback detail are configured separately from these dataset groups
(Section~\ref{sec:target-feedback}). They change the assistance provided during an episode, not
the sample's action model or reference depth.

\subsection{What a sample carries}
\label{sec:cp2seq-sample}

\begin{table}[h]
\caption{The files shipped with each CP2Seq sample. Only the first two, and only the final frame of
the second, are shown to the model.}
\label{tab:sample-files}
\begin{center}
\begin{tabular}{lp{0.68\textwidth}}
\multicolumn{1}{c}{\bf FILE}  &\multicolumn{1}{c}{\bf CONTENTS} \\
\hline \\
\texttt{cp.fold}    & The crease pattern, planarised. \textbf{The input} \\
\texttt{steps.fold} & The pattern plus the folded state after each step, as one multi-frame file. The final frame is \textbf{the other input}; the intermediate frames are held back \\
\texttt{seq.json}   & Every fold: line, direction, selected layers, creases made, coupling. \textbf{The ground truth}, never shown to the model \\
\texttt{meta.json}  & Difficulty metrics, degeneracy flags and generator version \\
\end{tabular}
\end{center}
\end{table}


\subsection{Two limits that belong in the paper rather than an appendix. --------MISSING}
\label{sec:cp2seq-limits}

\textbf{Caveat.} The first is an empty cell. Long sequences at low coupling are almost unreachable,
yielding 2% fact:corpus.synth.lLocalHits
hits in 16,000 attempts, because every all-layers fold thickens the stack, so a long sequence cannot
keep cutting few layers. Real folders reach that corner by \textbf{pre-creasing}, an operation this
action space excludes. The longer a real model runs, the more of it sits outside what the benchmark
can express.

\textbf{Caveat.} The second is depth. The some-layers action model stops at six folds in this release, which limits the sequence lengths represented for that action model. Longer reference sequences
require additional generation and evaluation before claims about their difficulty can be made.

% ===========================================================================
% FIGURE 3 -- THE DIFFICULTY GRID.  [TO DRAW]
% Depth on one axis, coupling on the other, one cell per stratum, shaded by how
% many samples landed there. The long-sequence low-coupling corner must be
% visibly EMPTY, and the caption names pre-creasing as the reason.
% ===========================================================================
\begin{figure}[h]
\draftimage{7}
\begin{center}\textit{\footnotesize Figure 3 placeholder: the
difficulty grid. Fold depth against coupling, shaded by sample count, with the long-sequence
low-coupling cell visibly empty.}\end{center}
\caption{\textbf{[DRAFT IMAGE]} The difficulty grid. The empty corner is long sequences at low coupling, which this action
space cannot reach because it excludes pre-creasing.}
\label{fig:grid}
\end{figure}

\subsection{Verifying the corpus itself}
\label{sec:cp2seq-verify}

Two checks run over every sample. Both replay the recorded sequence and compare the creases it
makes against the pattern stored beside it; what they do not share is the comparison. The first
groups edges into lines, merges collinear pieces and matches them within a tolerance. The second
was written after five false failures had come out of exactly that machinery, each one a defect
in the comparison rather than in the corpus, and reuses none of it. It merges each line's pieces
into maximal creased intervals, the quantity subdivision cannot change, and matches those one to
one. Comparing subdivided segments instead does not work: of 457 failures under that earlier
design, 364 were two sides carrying a different \emph{number} of segments, which is not a distance and
no tolerance can reach. The verdict is an absolute test at the radius at which the generator's own
planarizer identifies two points as one vertex, 1e-9, the floor of what the corpus records rather
than a constant tuned until the run passed. A depth-scaled ULP distance is reported alongside as a
diagnostic on how much of that margin is used; on a sampled subset the median match needed 16 ULP,
four orders inside the floor. Both checks pass on all 600% fact:corpus.release.total
samples of the release corpus.
